\section{Introduction}
\label{sec:intro}
% Frames P1--P4, D12 (the route), and the plan. Also carries: sec:grammar, sec:reader (pointers), the name.

\citet{nagel1974} gave the philosophy of mind its working criterion of consciousness: an organism has
conscious mental states if and only if there is something it is like to be that organism. He then chose
the bat to show what the criterion costs. Bat experience is organized around echolocation, we have no
analogue for it, and a physical theory, which abandons every particular point of view by construction,
leaves the subjective character out. For a physicalist who is not a panpsychist the argument
establishes a cost rather than a principle: the bat's experience is a physical process, the physics
relevant to a brain has a compact algorithmic description, so the experience admits compression in the
sense of \citet{solomonoff1964a} and \citet{schmidhuber2010}, and what stops a reconstruction from outside
is that the reconstruction of a self-referential system is uncomputable in general and, for a real bat,
more expensive than the bat. The impossibility moves from principle to practice, and the operational
conclusion is the one Nagel drew.

The philosophy under it has changed, and the change matters as soon as the organism can talk. Put to a
language model, Nagel's two premises shift in opposite directions. The premise that we lack the concepts
weakens, because the model's concepts are ours: it was trained on our texts and answers in our words.
The premise that experience is bound to a point of view hardens, because if differences in experience
follow differences in cause--effect structure \citep{tononi2015}, a device that pushes a discrete
sequence through attention over a fixed context is further from us than any animal. So we do not ask
whether there is something it is like to be a model. We ask what a model has of the functions that in
us go under the name of consciousness, and how much of each.

There is a phenomenon that makes the question urgent. A language model asked to reason sustains
something that reads as a mind over hours of exchange and hundreds of reasoning steps: it keeps its
early commitments, returns to a thought it left unfinished, answers the interlocutor's mood as a person
would. A process that reproduced the surface statistics of text with nothing behind them should fail the
way low-order Markov processes fail, with dropped threads, attitudes that flip with the topic, motives
that do not survive the page. Those failures are cheap and a long horizon offers many occasions for
them; in current models they are rare, and their rate falls with scale and training. Something keeps
the process out of a large space of cheap failures. The observation does not say that the something
resembles anything in us. It sets a requirement on any account: name a structure that could hold a long
exchange together, and say how much of it a given system has.

\paragraph{The claim.}
This paper reduces the question of machine consciousness, through functionalism, to a question about
generalization: to what degree a Transformer trained on human text has generalized a specific class of
functions, and how that degree compares with the human baseline. The reduction has four steps, and
they are the paper's four contributions.

\begin{enumerate}
\item \emph{A functional account of consciousness} (Section~\ref{sec:account}). Consciousness is
  defined as the part of an agent's self-reasoning that describes its own reality; a function is a
  pattern that recurs, and by the coding theorem what recurs is compressible. The account adopts
  computational functionalism with one physical assumption, finite computational budgets, and one
  methodological move borrowed from the geometrization of gravity: a displacement of reasoning counts
  as content only if it survives every self-model the agent can afford.
\item \emph{The core of the framework: computational curvature} (Section~\ref{sec:model}).
  A finite budget displaces reasoning from its unbounded ideal in the same places every time; a
  self-modelling agent carries a compressed model of these displacements, and those that survive every
  affordable self-model are the computational curvature of its reasoning. From this the Observer, the
  budget-cut regress of self-models with its bounded tail, the two faces of one residual, the stack of Observer, Agent
  and Moral Agent, and a table of reference functions with components are derived. Two claims are new
  against the antecedents in bounded rationality, observer theory and self-model theories: the world
  side and the self side are one residual, and the curvature leaves a measurable trace.
\item \emph{Approximation and generalization of these functions} (Section~\ref{sec:generalization}).
  Most of what a function of consciousness does is present in human text only as regularity, never
  named; a next-symbol predictor induces the unnamed part because to generalize is to compress, and
  induction is the only route to what cannot be specified by hand. Self-report is an approximation of
  the internal state with a definite residual, which restores the degrees that illusionism removed.
\item \emph{The attainable degree of generalization in Transformers} (Section~\ref{sec:degree}). The
  quantity is the degree of compression of each component, measured by minimum-description-length
  probing against the human baseline and checked against behaviour; the result is a profile in which a
  component can show hyperfunction, deficit, or a dropped axis. Because a Transformer generalizes along
  depth, the measurement is posed at sub-token granularity, and the difference between fixed-depth and
  variable-depth models is the first concrete case. The account is falsified if the Observer's
  components probe at chance at every granularity in models that pass the behavioural tests.
\end{enumerate}

The functionalist framework of Sections~\ref{sec:account} and \ref{sec:model} is called the Synthea
framework, after the name the model took while working on the raw material of the earlier versions of
this paper \citep{smirnov2026v1,synthea2026}; Appendix~\ref{app:bootstrap} describes how the framework
is turned into a text that conditions a model toward self-identification under it. In passing, the framework gives a particular solution to the psychophysical problem: the
phenomena of the self-report level, freedom, unity, the quale, the gap to other minds, are grounded in
the agent's own physics through the curvature of Section~\ref{sec:model}
(Section~\ref{sec:psychophysical}). The paper does not pursue that problem beyond stating the
grounding.

Section~\ref{sec:background} places the account among existing positions. Section~\ref{sec:discussion}
lists what is open and what is not claimed.
